%%%%%%%%%%%%%%%%%%%%%%%%%%%%%%%%%%%%%%%%%%%%%%%%%%%%%%%%%%%%%%%%%%%%%%
% Lớp 8 | Chương II. Hằng đẳng thức đáng nhớ và ứng dụng | Hiệu hai bình phương. Bình phương của một tổng hay một hiệu
% Dạng 8D2?6-2: Viết biểu thức dưới dạng bình phương, hiệu hai bình phương
%%%%%%%%%%%%%%%%%%%%%%%%%%%%%%%%%%%%%%%%%%%%%%%%%%%%%%%%%%%%%%%%%%%%%%


% nguồn: dan_tu_AI (nhập 14/11/2026)
\begin{bt}%[8D2H6-2]%[ma:f2591549]
	Viết các biểu thức sau dưới dạng bình phương của một tổng hoặc một hiệu:\\
	a) $x^2+6x+9$;\qquad b) $4y^2-4y+1$;\qquad c) $9a^2+12ab+4b^2$.
	\loigiai{
		a) $x^2+6x+9=x^2+2\cdot x\cdot 3+3^2=(x+3)^2$.\\
		b) $4y^2-4y+1=(2y)^2-2\cdot 2y\cdot 1+1^2=(2y-1)^2$.\\
		c) $9a^2+12ab+4b^2=(3a)^2+2\cdot 3a\cdot 2b+(2b)^2=(3a+2b)^2$.
	}
\end{bt}

% nguồn: dan_tu_AI (nhập 14/11/2026)
\begin{bt}%[8D2H6-2]%[ma:11e3b28c]
	Viết các biểu thức sau dưới dạng tích:\\
	a) $x^2-25$;\qquad b) $16y^2-9$;\qquad c) $(x+1)^2-4$.
	\loigiai{
		a) $x^2-25=x^2-5^2=(x-5)(x+5)$.\\
		b) $16y^2-9=(4y)^2-3^2=(4y-3)(4y+3)$.\\
		c) $(x+1)^2-4=(x+1)^2-2^2=(x+1-2)(x+1+2)=(x-1)(x+3)$.
	}
\end{bt}

% nguồn: dan_tu_AI (nhập 14/11/2026)
\begin{bt}%[8D2H6-2]%[ma:25aed345]
	Điền biểu thức thích hợp vào chỗ trống ($\ldots$) để được hằng đẳng thức đúng:\\
	a) $x^2+\ldots+16=(x+\ldots)^2$;\\
	b) $25a^2-10a+\ldots=(\ldots)^2$;\\
	c) $\ldots-12xy+9y^2=(2x-\ldots)^2$.
	\loigiai{
		a) $16=4^2$ nên $x^2+2\cdot x\cdot 4+4^2=x^2+8x+16=(x+4)^2$.\\
		b) $25a^2=(5a)^2$ và $10a=2\cdot 5a\cdot 1$ nên $25a^2-10a+1=(5a-1)^2$.\\
		c) $9y^2=(3y)^2$ và $12xy=2\cdot 2x\cdot 3y$ nên $4x^2-12xy+9y^2=(2x-3y)^2$.
	}
\end{bt}

% nguồn: dan_tu_AI (nhập 14/11/2026)
\begin{bt}%[8D2V6-2]%[ma:30a3e0fa]
	Viết các biểu thức sau dưới dạng hiệu hai bình phương rồi viết thành tích:\\
	a) $x^2+2x+1-y^2$;\\
	b) $(2x+3)^2-(x-1)^2$.
	\loigiai{
		a) $x^2+2x+1-y^2=(x+1)^2-y^2=(x+1-y)(x+1+y)$.\\
		b) $(2x+3)^2-(x-1)^2=\left[(2x+3)-(x-1)\right]\left[(2x+3)+(x-1)\right]$\\
		$=(2x+3-x+1)(2x+3+x-1)=(x+4)(3x+2)$.
	}
\end{bt}

% nguồn: dan_tu_AI (nhập 14/11/2026)
\begin{bt}%[8D2V6-2]%[ma:fe6d002a]
	a) Viết biểu thức $x^2-8x+20$ dưới dạng $(x-a)^2+b$ với $a$, $b$ là các số.\\
	b) Viết biểu thức $P=x^2-2xy+y^2+4x-4y+4$ dưới dạng bình phương của một tổng.
	\loigiai{
		a) $x^2-8x+20=x^2-2\cdot x\cdot 4+16+4=(x-4)^2+4$.\\
		Vậy $a=4$, $b=4$.\\
		b) Ta có $x^2-2xy+y^2=(x-y)^2$ và $4x-4y=4(x-y)$. Đặt $t=x-y$:
		\[ P=t^2+4t+4=t^2+2\cdot t\cdot 2+2^2=(t+2)^2. \]
		Vậy $P=(x-y+2)^2$.
	}
\end{bt}

% nguồn: dan_tu_AI (nhập 16/11/2026)
\begin{bt}%[8D2V6-2]%[ma:4d4f7d4b]
	Viết biểu thức $A = 9x^2 - 12x + 4 - 25y^2$ dưới dạng hiệu hai bình phương. Từ đó phân tích $A$ thành nhân tử.
	\loigiai{
		Nhóm các hạng tử để xuất hiện hằng đẳng thức bình phương của một hiệu:
		\begin{align*} 
			A &= (9x^2 - 12x + 4) - 25y^2 \\
			&= [(3x)^2 - 2 \cdot 3x \cdot 2 + 2^2] - (5y)^2 \\
			&= (3x-2)^2 - (5y)^2.
		\end{align*}
		Đến đây $A$ đã có dạng hiệu hai bình phương. Tiếp tục áp dụng hằng đẳng thức để phân tích thành nhân tử:
		\begin{align*} 
			A &= [(3x-2) - 5y][(3x-2) + 5y] \\
			&= (3x - 5y - 2)(3x + 5y - 2).
		\end{align*}
	}
\end{bt}
